\chapter{Scope, Novelty Boundary, and Anticipated Objections}\label{app:adversarial}

\section{Objection one: the circle is only complex inversion}
\textbf{Strong form of the objection.} Complex inversion maps a vertical line to a circle. Semicircular complex response plots are classical. Therefore the transport circle is a change of coordinates rather than a substantive result.

\textbf{Response.} The geometric fact is classical and is not claimed as new. The substantive result is the exact classification of the anchored chiral response condition in terms of the measured inverse response, together with the symmetry boundary and the explicit correlated counterexample. In the published chiral FHP kinetic theory, the circle is indeed the inverse image of a vertical shear eigenvalue line. Beyond that closure, the theorem identifies exactly what must remain vertical after all residues, projections, and correlation effects are included in the measured response.

\textbf{Boundary of the defense.} The theorem does not establish a new general theorem about Mobius transformations. Its physical value depends on whether the inverse response criterion clarifies and constrains the correlated chiral transport problem.

\section{Objection two: $C_6$ symmetry should already force the result}
\textbf{Strong form of the objection.} Hexagonal isotropy reduces the shear sector to two scalar coefficients, so the circle may be only a symmetry consequence.

\textbf{Response.} The commutant calculation proves the opposite. $C_6$ permits every matrix $aI+bJ$. Reflection makes $a$ even and $b$ odd in chirality displacement, but allows arbitrary even drift in $a$. The positive dissipation OU family chooses $a=a_0+\epsilon d^2$ and $b=cd$ and obeys the same symmetry form while violating the circle exactly. Symmetry explains the two coordinate complex structure, not dissipative flatness.

\textbf{Boundary of the defense.} A stronger microscopic symmetry or conservation law could still force $a$ to remain constant in a particular model. The thesis has not found or proved such a law for correlated chiral FHP.

\section{Objection three: the counterexample is not FHP}
\textbf{Strong form of the objection.} An Ornstein Uhlenbeck process is too simple and is not a cellular automaton, so it cannot decide whether the FHP circle survives correlations.

\textbf{Response.} It is not intended to decide the FHP outcome. It answers a different logical question: whether stationarity, positive dissipation, reflection covariance, and the two component rotational normal form are sufficient to protect the circle. The exact counterexample shows that they are not. Consequently any proof for chiral FHP must use additional model specific structure.

\textbf{Boundary of the defense.} The counterexample cannot be used as evidence that correlated FHP breaks the circle. It only proves that the broad structural assumptions do not determine the answer.

\section{Why this model rather than a direct Green Kubo simulation?}
A direct simulation would provide finite regulator evidence but would not by itself establish the structural necessity or sufficiency of circle preservation. The inverse response theorem answers that structural question exactly. Simulation remains valuable as a later model specific test of $G_{L,T}(p)$ or $G_L(\omega,p)$.

\section{Why the parity anchored circle?}
The parity symmetric point supplies a distinguished real response value because the Hall coordinate vanishes there under the stated reflection symmetry. The circle with diameter between zero and this anchor is exactly the image of the vertical inverse response line through $a_0$. Other circles can be considered, but they represent different geometric constraints and are not CRTC as defined here.

\section{Why not classify only by curvature?}
Curvature is local. The quartic drift witness has the same achiral curvature as the rigid circle but leaves it away from the anchor. A global classifier must therefore use the inverse response identity or an equivalent condition over the domain.

\section{Prior art challenge}
\textbf{Question.} Is this only the 2026 chiral FHP result rewritten in complex notation?

\textbf{Answer.} The published work supplies the kinetic coefficients and the chiral lattice gas setting \cite{Allocca2026}. The present structural theorem does not assume those coefficients. It applies to any nonzero measured complex response and remains valid if correlations renormalize both the effective kernel and the response residue. The counterexample and spectral overlap obstruction also lie outside the published kinetic derivation. Conversely, the thesis attributes the molecular chaos coefficients and does not claim their derivation as original.

\section{Overclaim audit of the abstract}
\begin{longtable}{p{.42\textwidth}p{.48\textwidth}}
\toprule
Abstract statement & Support and boundary \\
\midrule
\endhead
Deterministic endpoint reciprocity is exact. & Operator identities plus exhaustive finite implementation certificate and independent oracle. Does not imply a viscosity value.\\
Published Boltzmann coefficients lie on a circle. & Algebraic elimination of cited formulas plus exact rational tests. Conditional on the published molecular chaos theory.\\
Anchored circle membership is equivalent to constant $\RePart(1/Z)$. & Exact algebraic theorem for every nonzero response with stated anchor.\\
$C_6$ gives $aI+bJ$ but not constant $a$. & Exact commutant solution plus explicit symmetry allowed drift.\\
A positive dissipation correlated counterexample exists. & Exact OU Lyapunov and Green Kubo calculation plus rational grid test. Not an FHP surrogate.\\
Correlated FHP survival remains open. & No current repository proof establishes the required constant inverse real part under a correlated response prescription.\\
\bottomrule
\end{longtable}

\section{Defense questions}
The following questions should be answerable directly from the thesis and repository.
\begin{enumerate}[label=\arabic*.]
\item Where exactly does molecular chaos enter the circle result?
\item Which statement is proved by $C_6$ symmetry, and which statement is not?
\item Why is $\RePart(1/Z)$ preferred to a bare collision eigenvalue when a response residue can vary?
\item What role do long time tails play in defining the correlated question?
\item Why is local curvature insufficient for the global classification?
\item What additional data beyond eigenvalues are needed for a general Green Kubo resolvent matrix element?
\item Which result would have to be proved before Ward identity language becomes justified?
\item What, precisely, does the OU counterexample falsify and what does it not falsify?
\end{enumerate}
